\section{2022级工科数学分析下 A卷}

\subsection{资料来源}
\begin{itemize}
  \item 2022级A卷无答案.pdf（试题，扫描型 PDF，已按页面人工转写）。
  \item 2022级A卷学生答案.pdf（学生答案，扫描型 PDF；能辨认处用于校对，缺口处按教材方法补写解析）。
\end{itemize}

\subsection{试题}

\paragraph*{一、填空题（共 5 题，每题 2 分，共 10 分）}
\begin{enumerate}
  \item 微分方程 $y''+2y'+6y=0$ 的通解为 \underline{\hspace{4cm}}。
  \item 设函数 $u=2x^2+2y^2+3z^2$，则 $\operatorname{div}(\operatorname{grad}u)=$ \underline{\hspace{3cm}}。
  \item 设 $\Gamma$ 为 $y^2=2x$ 上从原点到 $(2,2)$ 的一段，则第一类曲线积分
  \[
    \int_\Gamma y\,ds=
  \]
  \underline{\hspace{3cm}}。
  \item 参数曲线
  \[
    \begin{cases}
      x=t-\cos t,\\
      y=\sin t,\\
      z=2t
    \end{cases}
  \]
  在 $t=0$ 对应的点处的切线方程为 \underline{\hspace{5cm}}。
  \item 设周期为 $2\pi$ 的函数
  \[
    f(x)=
    \begin{cases}
      -1, & -\pi<x\le 0,\\
      1+x, & 0<x\le \pi,
    \end{cases}
  \]
  则 $f(x)$ 的傅里叶级数在 $x=\pi$ 处收敛于 \underline{\hspace{3cm}}。
\end{enumerate}

\paragraph*{二、选择题（共 5 题，每题 2 分，共 10 分）}
\begin{enumerate}
  \item 关于未知函数 $y$ 的微分方程
  \[
    (y-\cos^2 x)\,dx+e^x\,dy=0
  \]
  是（\quad）。

  \begin{tabular}{llll}
    A. 可分离变量方程； & B. 一阶线性非齐次方程； &
    C. 一阶线性齐次方程； & D. 非线性方程。
  \end{tabular}

  \item 若二元函数 $f(x,y)$ 满足
  \[
    \lim_{(x,y)\to(0,0)}
    \frac{f(x,y)-f(0,0)+3x-5y}{\sqrt{x^2+y^2}}=0,
  \]
  则（\quad）。

  \begin{tabular}{ll}
    A. $df(0,0)=0$； & B. $df(0,0)=3dx-5dy$；\\
    C. $df(0,0)=-3dx+5dy$； & D. $df(0,0)$ 不存在。
  \end{tabular}

  \item 设函数 $f(x,y)$ 与 $\varphi(x,y)$ 均可微，且 $\varphi_y(x,y)\ne0$。若 $(x_0,y_0)$ 是 $f(x,y)$ 在约束条件
  $\varphi(x,y)=0$ 下的一个极值点，则下列选项正确的是（\quad）。

  \begin{tabular}{ll}
    A. 若 $f_x(x_0,y_0)=0$，则 $f_y(x_0,y_0)=0$； &
    B. 若 $f_x(x_0,y_0)=0$，则 $f_y(x_0,y_0)\ne0$；\\
    C. 若 $f_x(x_0,y_0)\ne0$，则 $f_y(x_0,y_0)=0$； &
    D. 若 $f_x(x_0,y_0)\ne0$，则 $f_y(x_0,y_0)\ne0$。
  \end{tabular}

  \item 函数 $\ln(1+x)$ 在 $x=0$ 处的泰勒展开式正确的是（\quad）。

  \begin{tabular}{ll}
    A. $\displaystyle \ln(1+x)=\sum_{n=1}^{\infty}\frac{(-1)^n}{n}x^n,\ x\in(-1,1]$； &
    B. $\displaystyle \ln(1+x)=\sum_{n=1}^{\infty}\frac{(-1)^{n-1}}{n}x^n,\ x\in(-1,1]$；\\
    C. $\displaystyle \ln(1+x)=-\sum_{n=1}^{\infty}\frac{1}{n}x^n,\ x\in[-1,1)$； &
    D. $\displaystyle \ln(1+x)=\sum_{n=1}^{\infty}\frac{1}{n}x^n,\ x\in[-1,1)$。
  \end{tabular}

  \item 使得级数
  \[
    \sum_{n=1}^{\infty}\frac{(-1)^n}{n^p}
  \]
  条件收敛的常数 $p$ 的取值范围是（\quad）。

  \begin{tabular}{llll}
    A. $p\le0$； & B. $0<p\le1$； & C. $0<p<1$； & D. $p>1$。
  \end{tabular}
\end{enumerate}

\paragraph*{三、计算题（共 3 题，每题 10 分，共 30 分）}
\begin{enumerate}
  \item 设
  \[
    z=\frac{1}{x}f(x+y,x-y)+g(xy),
  \]
  其中函数 $f$ 有二阶连续的偏导数，且 $g$ 二阶可导，计算 $\dfrac{\partial z}{\partial x}$ 和
  $\dfrac{\partial^2z}{\partial x\partial y}$。

  \item 计算累次积分
  \[
    \int_{-\sqrt2}^{0}dx\int_{-x}^{\sqrt{4-x^2}}\sqrt{x^2+y^2}\,dy
    +\int_{0}^{2}dx\int_{\sqrt{2x-x^2}}^{\sqrt{4-x^2}}\sqrt{x^2+y^2}\,dy.
  \]

  \item 计算锥面 $z=\sqrt{x^2+y^2}$ 被柱面 $x^2+y^2=4x$ 截下的部分锥面面积。
\end{enumerate}

\paragraph*{四、综合题（共 3 题，每题 10 分，共 30 分）}
\begin{enumerate}
  \item 设 $\Sigma$ 是
  \[
    \frac{x^2}{a^2}+\frac{y^2}{b^2}+\frac{z^2}{c^2}=1\quad (z\ge0)
  \]
  的上侧，且 $a,b,c$ 都是正实数。计算第二类曲面积分
  \[
    \iint_\Sigma (x^2-y\sin x)\,dy\,dz+(y^2-z^2)\,dz\,dx+(z^2-x^2)\,dx\,dy.
  \]

  \item 计算曲线积分
  \[
    \int_\Gamma
    \frac{(x+y)\,dx-(x-y)\,dy}{x^2+y^2},
  \]
  其中曲线 $\Gamma$ 是从点 $A(-1,0)$ 到点 $B(1,0)$ 的一条不经过原点的光滑曲线 $y=f(x)$，
  $-1\le x\le1$。

  \item 求幂级数
  \[
    \sum_{n=0}^{\infty}(n+1)x^{2n+1}
  \]
  的收敛域及和函数。
\end{enumerate}

\paragraph*{五、证明题（共 1 题，每题 10 分，共 10 分）}
证明函数项级数
\[
  \sum_{n=1}^{\infty}\sin\frac{n!x^n}{x^2+n^n}
\]
在区间 $[-2,2]$ 上一致收敛。

\paragraph*{六、应用题（共 1 题，每题 10 分，共 10 分）}
制作一个体积为 $2$ 的长方体盒子，利用拉格朗日乘数法求出长、宽和高分别为多少时才能使其表面积最小。

\subsection{答案与解析}

\paragraph*{一、填空题}
\begin{enumerate}
  \item 特征方程为 $r^2+2r+6=0$，根为 $r=-1\pm\sqrt5\,i$，故
  \[
    y=e^{-x}\left(C_1\cos\sqrt5\,x+C_2\sin\sqrt5\,x\right).
  \]
  \item $\operatorname{div}(\operatorname{grad}u)=\Delta u=u_{xx}+u_{yy}+u_{zz}=4+4+6=14$。
  \item 令 $y=t,\ x=\dfrac{t^2}{2}$，$0\le t\le2$，则 $ds=\sqrt{1+t^2}\,dt$，所以
  \[
    \int_\Gamma y\,ds=\int_0^2 t\sqrt{1+t^2}\,dt
    =\frac{(1+t^2)^{3/2}}{3}\bigg|_0^2
    =\frac{5\sqrt5-1}{3}.
  \]
  \item 当 $t=0$ 时点为 $(-1,0,0)$，切向量为 $(1,1,2)$，故切线方程为
  \[
    \frac{x+1}{1}=\frac{y}{1}=\frac{z}{2}.
  \]
  \item 傅里叶级数在跳跃点收敛于左右极限的平均值。因
  \[
    f(\pi-0)=1+\pi,\qquad f(\pi+0)=f(-\pi+0)=-1,
  \]
  所以收敛值为 $\dfrac{\pi}{2}$。
\end{enumerate}

\paragraph*{二、选择题}
\begin{enumerate}
  \item 方程可化为 \(y'+e^{-x}y=e^{-x}\cos^2x\)，是一阶线性非齐次方程。选 B。
  \item 由可微定义，线性主部为 \(-3x+5y\)，故 \(df(0,0)=-3\,dx+5\,dy\)。选 C。
  \item 约束极值必要条件给出 \(f_x\varphi_y-f_y\varphi_x=0\)。因 \(\varphi_y\ne0\)，若 \(f_x\ne0\) 则必有 \(f_y\ne0\)。选 D。
  \item \(\displaystyle \ln(1+x)=\sum_{n=1}^{\infty}\frac{(-1)^{n-1}}{n}x^n\)，收敛区间为 $(-1,1]$。选 B。
  \item 交错 $p$ 级数在 $p>0$ 时收敛，在 $p>1$ 时绝对收敛，所以条件收敛范围为 $0<p\le1$。选 B。
\end{enumerate}

\paragraph*{三、计算题}
\begin{enumerate}
  \item 记 $u=x+y,\ v=x-y$，并把 $f_1,f_2,f_{11},f_{22}$ 均理解为在 $(u,v)$ 处的取值，把 $g',g''$
  理解为在 $xy$ 处的取值。先对 $x$ 求导：
  \[
    \frac{\partial z}{\partial x}
    =-\frac{f}{x^2}+\frac{f_1+f_2}{x}+y\,g'(xy).
  \]
  再对 $y$ 求导，利用 $u_y=1,\ v_y=-1$，得
  \[
    \frac{\partial^2z}{\partial x\partial y}
    =-\frac{f_1-f_2}{x^2}+\frac{f_{11}-f_{22}}{x}+g'(xy)+xy\,g''(xy).
  \]

  \item 第一部分区域为极坐标下
  \[
    \frac{\pi}{2}\le\theta\le\frac{3\pi}{4},\qquad 0\le r\le2,
  \]
  第二部分区域为
  \[
    0\le\theta\le\frac{\pi}{2},\qquad 2\cos\theta\le r\le2.
  \]
  因 $\sqrt{x^2+y^2}=r$，面积元为 $r\,dr\,d\theta$，所以原积分为
  \[
    \int_0^{\pi/2}\int_{2\cos\theta}^{2}r^2\,dr\,d\theta
    +\int_{\pi/2}^{3\pi/4}\int_0^2r^2\,dr\,d\theta
    =2\pi-\frac{16}{9}.
  \]

  \item 锥面写成 $z=\sqrt{x^2+y^2}$ 时，
  \[
    dS=\sqrt{1+z_x^2+z_y^2}\,dx\,dy=\sqrt2\,dx\,dy.
  \]
  投影区域为圆 $x^2+y^2\le4x$，即半径为 $2$ 的圆盘，面积为 $4\pi$。故锥面面积为
  \[
    S=\sqrt2\cdot4\pi=4\sqrt2\,\pi.
  \]
\end{enumerate}

\paragraph*{四、综合题}
\begin{enumerate}
  \item 记
  \[
    P=x^2-y\sin x,\qquad Q=y^2-z^2,\qquad R=z^2-x^2.
  \]
  用底面
  \[
    \Sigma_0:\ z=0,\quad \frac{x^2}{a^2}+\frac{y^2}{b^2}\le1
  \]
  补成上半椭球体 $\Omega$ 的外侧闭曲面，底面取向向下。由高斯公式，
  \[
    I_\Sigma+I_{\Sigma_0}
    =\iiint_\Omega (P_x+Q_y+R_z)\,dV
    =\iiint_\Omega (2x-y\cos x+2y+2z)\,dV.
  \]
  上半椭球关于 $x,y$ 对称，含 $x,y$ 的奇函数项积分为 $0$，故
  \[
    I_\Sigma+I_{\Sigma_0}=\iiint_\Omega 2z\,dV
    =\frac{\pi ab c^2}{2}.
  \]
  底面向下时
  \[
    I_{\Sigma_0}=\iint_{\frac{x^2}{a^2}+\frac{y^2}{b^2}\le1}x^2\,dx\,dy
    =\frac{\pi a^3b}{4}.
  \]
  因而
  \[
    I_\Sigma=\frac{\pi ab c^2}{2}-\frac{\pi a^3b}{4}.
  \]

  \item 积分可拆成
  \[
    \frac{(x+y)\,dx-(x-y)\,dy}{x^2+y^2}
    =\frac{x\,dx+y\,dy}{x^2+y^2}
    +\frac{y\,dx-x\,dy}{x^2+y^2}
    =d\ln r-d\theta.
  \]
  两端点到原点的距离都为 $1$，故 $d\ln r$ 的积分为 $0$。剩下的值由曲线绕过原点的方向决定。若
  $f(0)>0$，曲线从上半平面绕过原点，$\theta$ 可由 $\pi$ 连续变到 $0$，积分为 $\pi$；若
  $f(0)<0$，曲线从下半平面绕过原点，积分为 $-\pi$。因此
  \[
    \int_\Gamma
    \frac{(x+y)\,dx-(x-y)\,dy}{x^2+y^2}
    =
    \begin{cases}
      \pi, & f(0)>0,\\
      -\pi, & f(0)<0.
    \end{cases}
  \]

  \item 令 $t=x^2$，则
  \[
    \sum_{n=0}^{\infty}(n+1)x^{2n+1}
    =x\sum_{n=0}^{\infty}(n+1)t^n.
  \]
  当 $|t|<1$ 即 $|x|<1$ 时，
  \[
    \sum_{n=0}^{\infty}(n+1)t^n=\frac{1}{(1-t)^2},
  \]
  所以和函数为
  \[
    S(x)=\frac{x}{(1-x^2)^2},\qquad |x|<1.
  \]
  当 $x=\pm1$ 时通项不趋于 $0$，故发散。因此收敛域为 $(-1,1)$。
\end{enumerate}

\paragraph*{五、证明题}
对任意 $x\in[-2,2]$，有
\[
  \left|\sin\frac{n!x^n}{x^2+n^n}\right|
  \le
  \left|\frac{n!x^n}{x^2+n^n}\right|
  \le
  \frac{n!\,2^n}{n^n}.
\]
令 $M_n=\dfrac{n!\,2^n}{n^n}$，则
\[
  \frac{M_{n+1}}{M_n}
  =2\left(\frac{n}{n+1}\right)^n\longrightarrow \frac{2}{e}<1.
\]
故 $\sum M_n$ 收敛。由魏尔斯特拉斯判别法，原函数项级数在 $[-2,2]$ 上一致收敛。

\paragraph*{六、应用题}
设长方体长、宽、高分别为 $x,y,z>0$，体积约束为 $xyz=2$，闭盒子的表面积为
\[
  S=2(xy+xz+yz).
\]
设
\[
  L=2(xy+xz+yz)+\lambda(xyz-2).
\]
由拉格朗日乘数法，
\[
  2(y+z)+\lambda yz=0,\qquad
  2(x+z)+\lambda xz=0,\qquad
  2(x+y)+\lambda xy=0.
\]
两两相减得 $x=y=z$。再由 $xyz=2$，得
\[
  x=y=z=\sqrt[3]{2}.
\]
因此长、宽和高都等于 $\sqrt[3]{2}$ 时表面积最小。
